\chapter*{Referências}
\addcontentsline{toc}{chapter}{Referências}
\begin{SingleSpacing}\raggedright
\noindent ASSOCIAÇÃO BRASILEIRA DE NORMAS TÉCNICAS (ABNT). \textbf{NBR 12244}: poço tubular: construção de poço tubular para captação de água subterrânea. Rio de Janeiro: ABNT, 2006.\par\medskip
\noindent ASSOCIAÇÃO BRASILEIRA DE NORMAS TÉCNICAS (ABNT). \textbf{NBR 12212}: projeto de poço tubular para captação de água subterrânea: procedimento. Rio de Janeiro: ABNT, 2017.\par\medskip
\noindent ASSOCIAÇÃO BRASILEIRA DE NORMAS TÉCNICAS (ABNT). \textbf{NBR 17076}: projeto de sistema de tratamento de esgoto de menor porte: requisitos. Rio de Janeiro: ABNT, 2024.\par\medskip
\noindent BAHIA. Governo do Estado. \textbf{Centrais de abastecimento de associações são referência de política pública de saneamento rural}. Salvador, 2022. Disponível em: \url{https://www.ba.gov.br/comunicacao/2022/05/noticias/centrais-de-abastecimento-de-associacoes-sao-referencia-de-politica-publica-de-saneamento-rural}.\par\medskip
\noindent BANCO MUNDIAL. \textbf{Quadro Ambiental e Social}. Washington, DC: Banco Mundial, 2017.\par\medskip
\noindent BRASIL. \textbf{Lei nº~9.433, de 8 de janeiro de 1997}. Institui a Política Nacional de Recursos Hídricos, cria o Sistema Nacional de Gerenciamento de Recursos Hídricos, regulamenta o inciso XIX do art. 21 da Constituição Federal, e altera o art. 1º da Lei nº~8.001, de 13 de março de 1990, que modificou a Lei nº~7.990, de 28 de dezembro de 1989. Brasília, DF: Presidência da República, 1997.\par\medskip
\noindent BRASIL. \textbf{Lei nº~11.107, de 6 de abril de 2005}. Dispõe sobre normas gerais de contratação de consórcios públicos e dá outras providências. Brasília, DF: Presidência da República, 2005.\par\medskip
\noindent BRASIL. \textbf{Lei nº~11.445, de 5 de janeiro de 2007}. Estabelece as diretrizes nacionais para o saneamento básico [...]. Brasília, DF: Presidência da República, 2007.\par\medskip
\noindent BRASIL. \textbf{Lei nº~13.019, de 31 de julho de 2014}. Estabelece o regime jurídico das parcerias entre a administração pública e as organizações da sociedade civil. Brasília, DF: Presidência da República, 2014.\par\medskip
\noindent BRASIL. \textbf{Lei nº~13.709, de 14 de agosto de 2018}. Lei Geral de Proteção de Dados Pessoais (LGPD). Brasília, DF: Presidência da República, 2018.\par\medskip
\noindent BRASIL. Ministério da Saúde. Fundação Nacional de Saúde. \textbf{Programa Nacional de Saneamento Rural}. Brasília, DF: Funasa, 2019.\par\medskip
\noindent BRASIL. \textbf{Lei nº~14.026, de 15 de julho de 2020}. Atualiza o marco legal do saneamento básico e altera a Lei nº~9.984, de 17 de julho de 2000, [...]. Brasília, DF: Presidência da República, 2020.\par\medskip
\noindent BRASIL. Ministério da Saúde. \textbf{Portaria GM/MS nº~888, de 4 de maio de 2021}. Altera o Anexo XX da Portaria de Consolidação GM/MS nº~5, de 28 de setembro de 2017, para dispor sobre os procedimentos de controle e de vigilância da qualidade da água para consumo humano e seu padrão de potabilidade. Brasília, DF: Ministério da Saúde, 2021.\par\medskip
\noindent BRASIL. \textbf{Decreto nº~11.599, de 12 de julho de 2023}. Dispõe sobre a prestação regionalizada dos serviços públicos de saneamento básico, o apoio técnico e financeiro de que trata o art. 13 da Lei nº~14.026, de 15 de julho de 2020, a alocação de recursos públicos federais e os financiamentos com recursos da União ou com recursos geridos ou operados por órgãos ou entidades da União de que trata o art. 50 da Lei nº~11.445, de 5 de janeiro de 2007. Brasília, DF: Presidência da República, 2023.\par\medskip
\noindent BRASIL. Ministério da Saúde. Fundação Nacional de Saúde. \textbf{Portaria nº~2.777, de 30 de julho de 2025}. Brasília, DF: Funasa, 2025. Disponível em: \url{https://bvs.saude.gov.br/bvs/saudelegis/funasa/2025/prt2777_05_08_2025.html}.\par\medskip
\noindent CEARÁ. Governo do Estado. \textbf{Modelo cearense de saneamento rural sustentável alcança mais de 1,1 milhão de pessoas no campo}. Fortaleza, 16 out. 2025. Disponível em: \url{https://www.ceara.gov.br/2025/10/16/modelo-cearense-de-saneamento-rural-sustentavel-alcanca-mais-de-11-milhao-de-pessoas-no-campo/}.\par\medskip
\noindent INSTITUTO ÁGUA E TERRA (IAT). \textbf{Nota técnica}: Componente 3 – Segurança Hídrica para o Saneamento; Subcomponente 3.2 – Saneamento Rural. Curitiba: IAT, 2026a. Documento interno, versão de 4 mar. 2026.\par\medskip
\noindent INSTITUTO ÁGUA E TERRA (IAT). \textbf{Nota técnica}: Subcomponente 3.2.b – implantação e reabilitação de sistemas coletivos de abastecimento de água e esgotamento sanitário no meio rural. Curitiba: IAT, 2026b. Documento interno.\par\medskip
\noindent INSTITUTO ÁGUA E TERRA (IAT); INSTITUTO DE DESENVOLVIMENTO RURAL DO PARANÁ (IDR-PARANÁ). \textbf{Saneamento rural}: áreas prioritárias, associações de municípios, CISPAR e concentração de domicílios rurais. Curitiba, 2025. 1 mapa, formato A3.\par\medskip
\noindent INSTITUTO BRASILEIRO DE GEOGRAFIA E ESTATÍSTICA (IBGE). \textbf{Censo Demográfico 2022}: agregados por setores censitários. Rio de Janeiro: IBGE, 2024.\par\medskip
\noindent PARANÁ. \textbf{Lei nº~12.726, de 26 de novembro de 1999}. Institui a Política Estadual de Recursos Hídricos e adota outras providências. Curitiba: Assembleia Legislativa do Estado do Paraná, 1999.\par\medskip
\noindent PARANÁ. Secretaria de Estado do Planejamento. \textbf{Planilha de custos do PSH-PR}: PSH\_CUSTO\_Consolidada, versão v2408. Curitiba: SEPL, 2026. Documento interno.\par\medskip
\noindent SANEPAR – COMPANHIA DE SANEAMENTO DO PARANÁ. \textbf{Manual do Programa Sanepar Rural}. Curitiba: Sanepar, [s.d.]. Documento interno.\par\medskip
\noindent TONETTI, Adriano Luiz \textit{et al.} \textbf{Tratamento de esgotos domésticos em comunidades isoladas}: referencial para a escolha de soluções. Campinas: Unicamp, 2018.\par\medskip
\end{SingleSpacing}
